\section{Problema General}
La Defensoría de Derechos Politécnicos (DDP) tiene como misión "proteger y garantizar los derechos politécnicos, además de promover el conocimiento, enseñanza y difusión de la cultura de los Derechos Humanos entre la comunidad politécnica" \cite{3}, sin embargo, se enfrentan a problemáticas dentro del organismo. Estas problemáticas ralentizan el proceso de resolución de una queja, pues es un departamento que maneja un gran volumen de quejas día a día.

Tras conocer su proceso, se identificaron problemas de organización (gran carga de trabajo, falta de resultados tempranos y carencia de trazabilidad hacia el quejoso), deficiencias en el manejo de expedientes ya que la falta de recursos y personal limita la agilidad del trámite y dificultades directamente con el quejoso, pues al no conocer el proceso ni los requerimientos, las quejas carecen de datos o información importante.

Todas estas problemáticas derivaron en la necesidad de mejorar el proceso de una queja, pues las problemáticas mencionadas en el párrafo anterior llevaron a la ralentización de los tiempos al momento de las resoluciones, siendo importante resolver las quejas en el menor tiempo posible, algo que es casi imposible para el reducido grupo de personas encargado de gestionar cientos de quejas.